\documentclass[11pt,letterpaper]{article}

% ------------------ Codificación e idioma ------------------
\usepackage[utf8]{inputenc}
\usepackage[T1]{fontenc}
\usepackage[spanish,es-noquoting,es-tabla]{babel}
\usepackage{lmodern}

% ------------------ Matemáticas ------------------
\usepackage{amsmath,amssymb}
\usepackage{siunitx}
\sisetup{output-decimal-marker={,}}

% ------------------ Geometría y estilo ------------------
\usepackage[margin=2.4cm,top=2.6cm,bottom=2.4cm]{geometry}
\usepackage{xcolor}
\usepackage{fancyhdr}
\usepackage{titlesec}
\usepackage{enumitem}
\usepackage{graphicx}
\usepackage{tcolorbox}
\usepackage{array}
\usepackage{tikz}
\usetikzlibrary{babel,arrows.meta}

\definecolor{UnalGreen}{RGB}{0,153,51}

\titleformat{\section}
{\normalfont\large\bfseries\color{UnalGreen}}{\thesection}{0.8em}{}
\titleformat{\subsection}
{\normalfont\normalsize\bfseries\color{UnalGreen!75!black}}{\thesubsection}{0.7em}{}
\titlespacing*{\section}{0pt}{1.2\baselineskip}{0.4\baselineskip}

\newcommand{\rulesec}{\vspace{-0.7\baselineskip}{\color{UnalGreen}\rule{\textwidth}{0.8pt}}\vspace{0.2\baselineskip}}

\pagestyle{fancy}
\fancyhf{}
\lhead{\small Fundamentos de Mecánica (1000019) -- Taller en clase}
\rhead{\small Semestre 2026-2}
\cfoot{\thepage}
\renewcommand{\headrulewidth}{0.4pt}

% ------------------ Macros ------------------
\newcommand{\vv}[1]{\vec{#1}}
\newcommand{\ii}{\hat{\imath}}
\newcommand{\jj}{\hat{\jmath}}
\newcommand{\kk}{\hat{k}}
\newcommand{\pto}[4]{$#1$(#2,\,#3,\,#4)}

\setlist[enumerate,1]{leftmargin=*,label=\textbf{\arabic*.},itemsep=0.45em}
\setlist[enumerate,2]{label=\textit{\alph*}),itemsep=0.15em,topsep=0.2em}

\begin{document}
	
	% ================= Portada =================
	\begin{center}
		{\Large\bfseries Taller en Clase: Magnitudes, Vectores y Cinemática}\\[0.3em]
		{\large\bfseries Grupo 14 -- Fundamentos de Mecánica (1000019)}\\[0.6em]
		{\small Trabajo en equipos de \textbf{4} personas $\;\cdot\;$ Duración: 1 hora}
	\end{center}
	
	\vspace{0.4em}
	
	\begin{tcolorbox}[colback=UnalGreen!4,colframe=UnalGreen!70!black,boxrule=0.6pt,
		left=6pt,right=6pt,top=5pt,bottom=5pt]
		\textbf{Reglas de trabajo}
		\begin{itemize}[leftmargin=1.2em,itemsep=0.15em,topsep=0.25em]
			\item Se entrega \textbf{una sola solución por equipo}, con los cuatro nombres.
			\item Cada resultado debe llevar \textbf{unidades} y un número razonable de cifras significativas.
			\item Antes de entregar cada parte, el equipo debe verificar la \textbf{consistencia dimensional} y el \textbf{orden de magnitud} de sus respuestas.
			\item Roles (\textbf{rotan en cada parte}): \textsc{coordinador} organiza y reparte; \textsc{calculista} ejecuta las cuentas; \textsc{escéptico} revisa unidades, signos y plausibilidad; \textsc{relator} escribe la solución final.
		\end{itemize}
	\end{tcolorbox}
	
	\vspace{0.3em}
	
	\noindent\begin{tabular}{@{}p{0.13\textwidth}p{0.42\textwidth}p{0.40\textwidth}@{}}
		\textbf{Equipo:} & \hrulefill & \\[0.8em]
		\textbf{Integrantes:} & \hrulefill & \hrulefill \\[0.8em]
		& \hrulefill & \hrulefill \\
	\end{tabular}
	
	% ================= Parte 1 =================
	\section{Unidades, dimensiones y estimación \normalfont\small\textit{}}
	\rulesec
	
	\begin{enumerate}
		\item La rapidez del sonido en aire a \SI{20}{\celsius} es \SI{343}{m/s}.
		\begin{enumerate}
			\item Exprésela en km/h y en mi/h ($1\ \text{mi}=\SI{1609}{m}$).
			\item Un rayo cae y el trueno se escucha \SI{4,5}{s} después. ¿A qué distancia cayó el rayo? (La luz llega prácticamente de forma instantánea; justifiquen esta afirmación con un cálculo).
		\end{enumerate}
		
		\item El radio de un núcleo de uranio es aproximadamente \SI{7,4e-15}{m} y el radio del átomo de hidrógeno es \SI{0,529}{\text{\AA}}.
		\begin{enumerate}
			\item Escriban ambas cantidades en fm y en nm.
			\item ¿Cuántas veces más grande es el átomo que el núcleo? Si el núcleo tuviera el tamaño de una canica (\SI{1}{cm} de radio), ¿cuál sería el radio del átomo en metros? Comparen ese resultado con una distancia conocida de Bogotá.
		\end{enumerate}
		
		\item La rapidez $v$ con la que se propaga una onda en una cuerda depende de la tensión $T$ (unidades de fuerza) y de la densidad lineal de masa $\mu$ (masa por unidad de longitud). Suponiendo que $v = k\,T^{a}\mu^{b}$, con $k$ adimensional:
		\begin{enumerate}
			\item Determinen $a$ y $b$ por análisis dimensional.
			\item ¿Puede el análisis dimensional determinar el valor de $k$? Justifiquen.
		\end{enumerate}
		
	\end{enumerate}
	
	% ================= Parte 2 =================
	\section{Vectores \normalfont\small\textit{}}
	\rulesec
	
	\begin{enumerate}[resume]
		\item Sean $\vv{A} = 4\ii - \jj + 2\kk$, $\vv{B} = -2\ii + 3\jj + \kk$ y $\vv{C} = \ii + \jj - 3\kk$. Determinen:
		\begin{enumerate}
			\item $\vv{A} + \vv{B} - \vv{C}$ \quad y \quad $2\vv{A} - 3\vv{B} + \vv{C}$
			\item $|\vv{A}|$, $|\vv{B}|$ \ y \ $|\vv{A}+\vv{B}|$. ¿Es $|\vv{A}+\vv{B}| = |\vv{A}|+|\vv{B}|$?
			\item $\vv{A}\cdot\vv{B}$ \quad y \quad $\vv{A}\times\vv{B}$
			\item El vector unitario $\hat{u}$ en la dirección de $\vv{C}$ (verifiquen que $|\hat u| = 1$).
		\end{enumerate}
		\textbf{Verificación del escéptico:} comprueben explícitamente que $\vv{A}\times\vv{B}$ es perpendicular a $\vv{A}$ y a $\vv{B}$.
		
		
		\item \textbf{Preguntas conceptuales.} Respondan \textit{verdadero} o \textit{falso} y \textbf{justifiquen con un argumento o un contraejemplo}:
		\begin{enumerate}
			\item Si $|\vv{A}+\vv{B}| = |\vv{A}| + |\vv{B}|$, entonces $\vv{A}$ y $\vv{B}$ son paralelos.
			\item Si $\vv{A}\cdot\vv{B} = 0$, necesariamente $\vv{A}=0$ o $\vv{B}=0$.
			\item Siempre se cumple $|\vv{A}-\vv{B}| \le |\vv{A}| + |\vv{B}|$.
			\item $\vv{A}\times\vv{B} = \vv{B}\times\vv{A}$.
			\item Si $\vv{A}\cdot\vv{B} = \vv{A}\cdot\vv{C}$ con $\vv{A}\neq 0$, entonces $\vv{B} = \vv{C}$.
		\end{enumerate}
	\end{enumerate}
	
	\newpage
	
	% ================= Parte 3 =================
	\section{Cinemática en una dimensión \normalfont\small\textit{}}
	\rulesec
	
	\begin{enumerate}[resume]

		
		\item Un objeto se mueve a lo largo del eje $x$ con la velocidad $v_x(t)$ mostrada en la Figura~\ref{fig:vt}. En $t=0$ su posición es $x_0 = 0$.
		
		\begin{figure}[h!]
			\centering
			\begin{tikzpicture}[scale=1.0]
				% ejes
				\draw[->,thick] (-0.3,0) -- (10.8,0) node[below right] {$t$ (s)};
				\draw[->,thick] (0,-2.6) -- (0,2.9) node[left] {$v_x$ (m/s)};
				% grilla
				\draw[gray!30,very thin] (0,-2.5) grid[xstep=1,ystep=0.5] (10.5,2.5);
				% marcas x
				\foreach \x in {1,...,10} \draw (\x,0.08) -- (\x,-0.08);
				\foreach \x in {2,4,6,8,10} \node[below,font=\small] at (\x,-0.12) {\x};
				% marcas y
				\foreach \y/\l in {2/4,1/2,-1/-2,-2/-4} {\draw (0.08,\y) -- (-0.08,\y); \node[left,font=\small] at (-0.12,\y) {\l};}
				% curva: (0,0)->(2,4)->(4,4)->(8,-4)->(10,-4)
				\draw[UnalGreen,ultra thick] (0,0) -- (2,2) -- (4,2) -- (8,-2) -- (10,-2);
				\fill[UnalGreen] (0,0) circle (1.6pt) (2,2) circle (1.6pt) (4,2) circle (1.6pt) (8,-2) circle (1.6pt) (10,-2) circle (1.6pt);
			\end{tikzpicture}
			\caption{Velocidad del objeto en función del tiempo.}
			\label{fig:vt}
		\end{figure}
		
		\begin{enumerate}
			\item Determinen la aceleración en cada uno de los cuatro intervalos.
			\item Calculen la posición del objeto en $t = 4$ s, $t = 6$ s y $t = 10$ s.
			\item ¿En qué instante el objeto está \textbf{más alejado} del origen? ¿Cuál es esa distancia?
			\item Comparen el \textbf{desplazamiento} y la \textbf{distancia total recorrida} entre $t=0$ y $t=10$ s.
			\item Dibujen la gráfica de $a_x$ vs.\ $t$ en la cuadrícula siguiente.
		\end{enumerate}
		
		\begin{figure}[h!]
			\centering
			\begin{tikzpicture}[scale=0.95]
				\draw[gray!35,very thin] (0,-1.5) grid[xstep=0.5,ystep=0.5] (10.5,1.5);
				\draw[->,thick] (-0.3,0) -- (10.8,0) node[below right] {$t$ (s)};
				\draw[->,thick] (0,-1.8) -- (0,1.9) node[left] {$a_x$ (m/s$^2$)};
				\foreach \x in {1,...,10} \draw (\x,0.08) -- (\x,-0.08);
				\foreach \x in {2,4,6,8,10} \node[below,font=\small] at (\x,-0.12) {\x};
				\foreach \y in {1,-1} {\draw (0.08,\y) -- (-0.08,\y); \node[left,font=\small] at (-0.12,\y) {\pgfmathparse{int(\y*2)}\pgfmathresult};}
			\end{tikzpicture}
			\caption{Cuadrícula para la gráfica de $a_x$ vs.\ $t$.}
		\end{figure}
		
		\newpage
		
		\item Desde la azotea de un edificio se lanza una piedra \textbf{verticalmente hacia arriba} con rapidez inicial \SI{15}{m/s}. La piedra llega al suelo \SI{5,0}{s} después de ser lanzada. Tomen $g = \SI{9,8}{m/s^2}$ y desprecien la resistencia del aire.
		\begin{enumerate}
			\item Definan explícitamente su sistema de referencia (origen y sentido positivo) antes de plantear las ecuaciones.
			\item ¿Qué altura máxima alcanza la piedra por encima del punto de lanzamiento y en qué instante?
			\item ¿Cuál es la altura del edificio?
			\item ¿Con qué rapidez llega al suelo?
			\item ¿En qué instante(s) la piedra pasa por la altura del punto de lanzamiento? Interpreten \textbf{todas} las raíces de la ecuación cuadrática.
		\end{enumerate}
		
		
		\item \textbf{Preguntas conceptuales.} Respondan y justifiquen con un ejemplo concreto:
		\begin{enumerate}
			\item ¿Puede un objeto tener velocidad nula y aceleración distinta de cero?
			\item ¿Puede un objeto tener aceleración negativa y estar aumentando su rapidez?
			\item Si la gráfica $x$ vs.\ $t$ es una recta con pendiente negativa, ¿qué se puede afirmar sobre $v_x$ y $a_x$?
			\item Dos objetos parten del mismo punto al mismo tiempo y llegan al mismo punto en el mismo instante. ¿Tienen necesariamente la misma velocidad media? ¿La misma rapidez media?
		\end{enumerate}
	\end{enumerate}
	
	\newpage
	
	
	\vspace{0.6em}
	\noindent\small\textit{Ejercicios adaptados del material del curso y del taller de la profesora Edna Carolina Pinilla Beltrán. Trabajo en grupos de cuatro.}
	
\end{document}
